\section{自同构的行列式}

我们首先假设体 D 是交换的。鉴于上面的注，V 的自同构 u 的行列式是唯一的元素 \(\det u\) 属于 \(D^{*}\)，使得我们有

\[
\omega (u (v _ {1}), \ldots , u (v _ {n})) = (\det u) \omega (v _ {1}, \ldots , v _ {n})\tag{14}
\]

对 V 的每个基 \((v_{1},\ldots ,v_{n})\) 与每个元素 \(\omega\)（属于 \(\Omega (\mathrm{V})\)）成立。

现在回到不再假设 D 交换的情形。

命题 2. --- a) 设 u 是 V 的一个自同构。存在 \(D_{ab}^{*}\) 的唯一元素，记作 det u 并称为 u 的行列式，使得我们有

\[
\omega (u (v _ {1}), \ldots , u (v _ {n})) = (\det u) \omega (v _ {1}, \ldots , v _ {n})\tag{15}
\]

对 V 的每个基 \((v_{1},\ldots ,v_{n})\) 与每个 \(\omega\)（在 \(\Omega (\mathrm{V})\) 中）成立

\begin{enumerate}
\def\labelenumi{\alph{enumi})}
\setcounter{enumi}{1}
\tightlist
\item
  映射 \(u \mapsto \det u\)（从 \(\mathbf{GL}(\mathrm{V})\) 到 \(\mathrm{D}_{\mathrm{ab}}^{*}\)）是一个群同态。
\end{enumerate}

设 \(\omega_0\) 是 \(\Omega (\mathrm{V})\) 的一个元素。映射 \((v_{1},\ldots ,v_{n})\mapsto\) \(\omega_0(u(v_1),\dots ,u(v_n))\)（从 B(V) 到 \(\mathrm{D}_{\mathrm{ab}}^{*}\)）属于 \(\Omega (\mathrm{V})\)。由 VIII, 第 451 页的推论 2，存在唯一的元素 \(t\) 属于 \(\mathrm{D}_{\mathrm{ab}}^{*}\)，使得我们有

\[
\omega_ {0} (u (v _ {1}), \ldots , u (v _ {n})) = t \omega_ {0} (v _ {1}, \ldots , v _ {n})
\]

对 \((v_{1},\ldots ,v_{n})\)（在 \(\mathrm{B(V)}\) 中）成立。若 \(\omega\) 是 \(\Omega (\mathrm{V})\) 的另一个元素，则存在一个元素 \(s\) 属于 \(\mathrm{D}_{\mathrm{ab}}^{*}\)，使得

\[
\omega (v _ {1}, \ldots , v _ {n}) = s \omega_ {0} (v _ {1}, \ldots , v _ {n})
\]

对 V 的每个基 \((v_{1},\ldots ,v_{n})\) 成立（同前所引）。我们立即推出关系式

\[
\omega (u (v _ {1}), \dots , u (v _ {n})) = t \omega (v _ {1}, \dots , v _ {n}).
\]

这就证明了 a)；断言 b) 是 a) 的直接推论。

